% ACL / EACL / NAACL / EMNLP / ARR Conference Paper Template
% 
% This template demonstrates ACL Anthology style using natbib + acl_natbib.bst.
% 
% For ARR submissions: use \usepackage[review]{acl}
% For non-anonymous preprint: use \usepackage[nonaclfinalcopy]{acl}
% For final camera-ready: use \usepackage{acl} + \aclfinalcopy
% 
% Required style files (download from https://github.com/acl-org/acl-style-files):
%   - acl.sty
%   - acl_natbib.bst
% 
% Compilation:
%   pdflatex paper.tex
%   bibtex paper
%   pdflatex paper.tex
%   pdflatex paper.tex

\documentclass[11pt,a4paper]{article}
\usepackage[T1]{fontenc}

% --- Choose ONE of the following ---
% \usepackage[review]{acl}          % ARR review (double-blind, line numbers)
% \usepackage[nonaclfinalcopy]{acl} % Non-anonymous preprint
\usepackage{acl}                    % Final camera-ready (requires \aclfinalcopy below)

\usepackage{natbib}
\usepackage{amsmath,amssymb,amsfonts}
\usepackage{graphicx}
\usepackage{xcolor}
\usepackage{hyperref}
\usepackage{xspace}
\usepackage{url}
\usepackage{booktabs}
\usepackage{microtype}
\usepackage{float}

% Enable this for final camera-ready version
% \aclfinalcopy

% ============================================================================
% Title & Authors
% ============================================================================

\title{Your Paper Title: A Descriptive and Informative Title}

% For ARR review (double-blind): remove author block entirely
\author{
  First Author\textsuperscript{1},
  Second Author\textsuperscript{2},
  Third Author\textsuperscript{1}}
\affiliations{
  \textsuperscript{1}Department of Computer Science, University One\\
  \textsuperscript{2}School of Computing, University Two}
\email{\{first.author, third.author\}@univ1.edu, second.author@univ2.edu}

% ============================================================================
% Document
% ============================================================================

\begin{document}
\maketitle

% ============================================================================
\begin{abstract}
This is where your abstract goes. ACL venues typically expect 150--250 words. 
The abstract should concisely summarize:
(1) the problem or research question,
(2) your approach or methodology,
(3) key results or findings, and
(4) main conclusions or implications.

For ARR submissions, keep the abstract self-contained and avoid references 
that reveal author identity.
\end{abstract}

% ============================================================================
\section{Introduction}

The introduction should establish the context and motivation for your work. 
Start broad (the general problem area), then narrow to the specific problem 
you address.

Key elements of an ACL-style introduction:
\begin{itemize}
  \item \textbf{Motivation:} Why is this problem important? What gap exists 
    in current approaches? \citep{vaswani2017attention}
  \item \textbf{Problem statement:} Clearly state the research question or 
    objective.
  \item \textbf{Approach overview:} Briefly describe your method (1--2 sentences).
  \item \textbf{Contributions:} List 3--5 concrete contributions, typically 
    as a bulleted or numbered list.
  \item \textbf{Paper structure:} Optional roadmap sentence (``Section 2 
    reviews related work...'').
\end{itemize}

Our contributions are:
\begin{enumerate}
  \item First contribution item — be specific and measurable.
  \item Second contribution item — avoid vague claims.
  \item Third contribution item — state what is \emph{new}.
\end{enumerate}

% ============================================================================
\section{Related Work}

Organize related work thematically, not chronologically. For an NLP paper, 
common organizational themes include:

\subsection{Theme One: Background Area}

Discuss foundational and closely related work. Group papers by approach or 
contribution, not by paper. Compare and contrast rather than summarizing 
each paper individually.

\citet{devlin2019bert} introduced BERT, which...

\subsection{Theme Two: Alternative Approaches}

Discuss competing approaches and explain how your work differs.

\citet{brown2020language} demonstrated that...

\subsection{Our Position}

Explicitly state how your work differs from and improves upon existing work. 
This subsection is common in ACL papers and helps reviewers understand your 
novelty.

% ============================================================================
\section{Methodology / Approach}

Describe your method in sufficient detail that a knowledgeable reader could 
reproduce it. For NLP papers, this typically includes:

\subsection{Problem Formulation}

Formally define the task. Use mathematical notation where appropriate.

Given an input sequence $x = (x_1, x_2, \ldots, x_n)$, the goal is to...

\subsection{Model Architecture}

Describe your model. Use diagrams (with \texttt{figure} environments) for 
complex architectures.

\begin{figure}[t]
  \centering
  % \includegraphics[width=0.8\columnwidth]{figures/architecture.pdf}
  \fbox{\parbox{0.6\columnwidth}{\centering\vspace{3cm}Architecture diagram\\placeholder}}
  \caption{Overview of our proposed model architecture. 
    (Replace with actual figure.)}
  \label{fig:architecture}
\end{figure}

Our model consists of three main components:
\begin{enumerate}
  \item \textbf{Encoder:} A transformer-based encoder that...
  \item \textbf{Interaction module:} A cross-attention mechanism that...
  \item \textbf{Decoder/classifier:} A linear projection that...
\end{enumerate}

\subsection{Training}

Describe training details: dataset(s), hyperparameters, optimization, 
computational resources.

\begin{table}[t]
  \centering
  \begin{tabular}{@{}ll@{}}
    \toprule
    \textbf{Hyperparameter} & \textbf{Value} \\
    \midrule
    Learning rate & $5 \times 10^{-5}$ \\
    Batch size & 32 \\
    Epochs & 3 \\
    Optimizer & AdamW \\
    Warmup steps & 500 \\
    Max sequence length & 512 \\
    \bottomrule
  \end{tabular}
  \caption{Training hyperparameters.}
  \label{tab:hyperparams}
\end{table}

% ============================================================================
\section{Experimental Setup}

\subsection{Datasets}

Describe each dataset: size, source, splits, preprocessing. Use a table for 
multiple datasets.

\begin{table}[t]
  \centering
  \begin{tabular}{@{}lrrr@{}}
    \toprule
    \textbf{Dataset} & \textbf{Train} & \textbf{Dev} & \textbf{Test} \\
    \midrule
    Dataset A       & 10,000 & 1,000 & 1,000 \\
    Dataset B       & 50,000 & 5,000 & 5,000 \\
    \bottomrule
  \end{tabular}
  \caption{Dataset statistics.}
  \label{tab:datasets}
\end{table}

\subsection{Baselines}

List and briefly describe each baseline. Distinguish between:
\begin{itemize}
  \item \textbf{Standard baselines:} Well-known approaches (e.g., BERT-base, 
    RoBERTa-large).
  \item \textbf{State-of-the-art:} Current best published results.
  \item \textbf{Ablations:} Variants of your model for analysis.
\end{itemize}

\subsection{Evaluation Metrics}

Define all metrics used. For NLP tasks, common metrics include:
\begin{itemize}
  \item Accuracy, F1 (classification)
  \item BLEU, ROUGE, METEOR (generation)
  \item Exact Match, F1 (extractive QA)
  \item BERTScore, BLEURT (semantic similarity)
\end{itemize}

% ============================================================================
\section{Results}

\subsection{Main Results}

Present your primary results. Use tables with clear formatting. Bold the best 
score in each column. Report statistical significance where appropriate.

\begin{table}[t]
  \centering
  \begin{tabular}{@{}lcccc@{}}
    \toprule
    \textbf{Model} & \textbf{Dataset A} & \textbf{Dataset B} & 
    \textbf{Dataset C} & \textbf{Avg} \\
    \midrule
    Baseline 1 & 72.3 & 68.1 & 75.4 & 71.9 \\
    Baseline 2 & 74.1 & 70.2 & 76.8 & 73.7 \\
    SOTA       & 76.5 & 72.8 & 78.9 & 76.1 \\
    \midrule
    Ours       & \textbf{78.2} & \textbf{74.5} & \textbf{80.1} & 
    \textbf{77.6} \\
    \bottomrule
  \end{tabular}
  \caption{Main results. Best scores in \textbf{bold}. 
    Our model outperforms all baselines across all datasets.}
  \label{tab:main-results}
\end{table}

\subsection{Ablation Studies}

Analyze the contribution of each component. Remove or modify individual 
components and measure the impact.

\begin{table}[t]
  \centering
  \begin{tabular}{@{}lc@{}}
    \toprule
    \textbf{Configuration} & \textbf{F1} \\
    \midrule
    Full model              & \textbf{78.2} \\
    -- Component A          & 75.1 \\
    -- Component B          & 76.3 \\
    -- Component A \& B     & 72.8 \\
    With smaller model      & 74.9 \\
    With less training data & 73.5 \\
    \bottomrule
  \end{tabular}
  \caption{Ablation results. Removing any component degrades performance.}
  \label{tab:ablations}
\end{table}

\subsection{Analysis and Discussion}

Go beyond numbers. Analyze \emph{why} your model works:
\begin{itemize}
  \item Qualitative examples (success and failure cases)
  \item Error analysis (what types of errors does your model still make?)
  \item Attention visualization or feature analysis
  \item Comparison with human performance
\end{itemize}

% ============================================================================
\section*{Limitations}

ACL venues strongly encourage (and some require) a Limitations section. 
Be honest about:
\begin{itemize}
  \item \textbf{Scope:} What tasks, languages, or domains does your approach 
    NOT cover?
  \item \textbf{Data:} Biases in training data, limited language coverage, 
    dataset artifacts.
  \item \textbf{Computational cost:} Training time, inference latency, model 
    size.
  \item \textbf{Assumptions:} Simplifying assumptions that may not hold in 
    practice.
  \item \textbf{Evaluation:} Limits of automatic metrics, lack of human 
    evaluation.
  \item \textbf{Generalizability:} Results on specific benchmarks may not 
    transfer.
\end{itemize}

% ============================================================================
\section*{Ethical Considerations / Broader Impact}

Some ACL venues require an ethics statement. Consider:
\begin{itemize}
  \item Potential misuse of the technology
  \item Bias and fairness implications
  \item Environmental impact (carbon footprint of training)
  \item Privacy concerns (if using personal data)
  \item Intended use and users
\end{itemize}

% ============================================================================
\section*{Acknowledgments}

% For ARR review: comment out or remove this section
We thank the anonymous reviewers for their helpful feedback. 
This work was supported by [Grant/Agency].

% ============================================================================
% Bibliography
% ============================================================================

\bibliographystyle{acl_natbib}
\bibliography{references}

% ============================================================================
% Appendix (optional, usually not included in page limit)
% ============================================================================

% \appendix
% \section{Additional Experiments}
% \section{Implementation Details}
% \section{Full Results}

\end{document}
